\documentclass{article}


\usepackage[preprint]{utils_packages/neurips_2023}



\usepackage[utf8]{inputenc} % allow utf-8 input
\usepackage[T1]{fontenc}    % use 8-bit T1 fonts
\usepackage{hyperref}       % hyperlinks
\usepackage{url}            % simple URL typesetting
\usepackage{booktabs}       % professional-quality tables
\usepackage{amsfonts}       % blackboard math symbols
\usepackage{nicefrac}       % compact symbols for 1/2, etc.
\usepackage{microtype}      % microtypography
\usepackage{xcolor}         % colorsp
\usepackage{float}
\usepackage{listings}

\urlstyle{same}
\hypersetup{
  hidelinks,
  pdfcreator={LaTeX via pandoc}}

\author{}
\date{}


\newcommand{\Eric}[1]{{\color{mint}[Eric: #1]}}
\definecolor{mint}{RGB}{62, 180, 137}
% [Custom packages] from here down 
% [Custom packages] from here down 
\usepackage{float}
\usepackage{todonotes}
\usepackage{amsmath}
\usepackage{algorithm}
\usepackage{algpseudocode}
\usepackage{enumitem}
\usepackage{caption}
\usepackage{subcaption}
% \usepackage{biblatex}
% \addbibresource{refs.bib}

\newcommand{\x}{\boldsymbol{x}}
% \newcommand{\y}{\boldsymbol{y}}
\newcommand{\X}{\boldsymbol{X}}
\newcommand{\w}{\boldsymbol{w}}
\newcommand{\sign}{\text{sign}}
\newcommand{\ub}{\boldsymbol{u}}
\newcommand{\z}{\boldsymbol{z}}
\newcommand{\vb}{\boldsymbol{v}}

\newcommand{\numberset}{\mathbb}
\newcommand{\N}{\numberset{N}}
\newcommand{\Z}{\numberset{Z}}
\newcommand{\Q}{\numberset{Q}}
\newcommand{\R}{\numberset{R}}
\newcommand{\C}{\numberset{C}}
\newcommand{\E}{\numberset{E}}
\newcommand{\Prob}{\numberset{P}} 

\usepackage{mathtools}
\DeclarePairedDelimiter\ceil{\lceil}{\rceil}
\DeclarePairedDelimiter\floor{\lfloor}{\rfloor}
\newcommand{\norm}[1]{\left\lVert#1\right\rVert}

\usepackage{nomencl}
\def\bmat#1{\left[\begin{array}{#1}}
\def\emat{\end{array}\right]}
\def\bv {\bmat{c} }
\def\ev {\emat}

\newcommand {\bsis} {\left\{ \begin{array} }
\newcommand {\esis} {\end{array}\right.}

\def\Sum#1#2{\sum\limits_{#1}^{#2}}
\def\set#1#2{\{ \; #1 \;:\;#2\;\}} % creates a set: { #1 : #2 }



\newcommand{\y}{\hat{y}}

\newcommand{\tauB}{\boldsymbol{\tau}}
\newcommand{\xiB}{\boldsymbol{\xi}}
\newcommand{\thetaB}{\boldsymbol{\theta}}
\newcommand{\thB}{{\boldsymbol{\theta}}}
\newcommand{\alB}{{\boldsymbol{\alpha}}}
\newcommand{\betB}{{\boldsymbol{\beta}}}
\newcommand{\yB}{{\boldsymbol{y}}}


\setcounter{MaxMatrixCols}{20}
\newtheorem{remark}{Remark}[section]


% [Custom packages] from here up 

\usepackage{amsmath}

\usepackage{comment}

% [Custom packages] from here up 


\title{EvoScale: Distributed Evolutionary Search\\
  with Online Reinforcement Learning}

% The \author macro works with any number of authors. There are two commands
% used to separate the names and addresses of multiple authors: \And and \AND.
%
% Using \And between authors leaves it to LaTeX to determine where to break the
% lines. Using \AND forces a line break at that point. So, if LaTeX puts 3 of 4
% authors names on the first line, and the last on the second line, try using
% \AND instead of \And before the third author name.
\author{
Group 7 \\ Matteo Guarrera \quad Elizabeth Polito \quad Harper Lyu \\ Mert Cemri \quad Sultan Daniels \quad Eric Xu
}

\begin{document}
\maketitle
% \vspace{-5cm}
\author{}
% 
\vspace{-0.5cm}

\vspace{-0.3cm}



\section*{Introduction}

LLM-driven evolutionary search has recently emerged as a powerful paradigm
for automated program optimization. The system maintains a population of
candidate programs and, at each iteration, prompts an LLM with a
high-scoring parent and several inspirations to propose an improvement.
AdaEvolve~\cite{adaevolve2025} is the current state of the art: it
introduces adaptive multi-island search that dynamically modulates
exploration intensity per island via an accumulated improvement signal and
routes compute across islands with a UCB bandit, requiring no per-task
configuration. The results are strong: approximately 50\% solve rate on
ARC-AGI-2, 3$\times$ mean performance over single-call baselines on 172
Frontier-CS algorithm design tasks, and competitive results against
AlphaEvolve on mathematical optimization.

Despite these results, the framework has two structural limitations. First,
execution is fully sequential: each iteration is one LLM call followed by
one evaluation. Evaluation of compiled or executable programs can take tens
of seconds, during which the GPU sits idle. Second, the LLM is never
updated. Every run generates thousands of structured (state, mutation,
outcome) tuples and discards all of them. The mutation operator is fixed
regardless of how much compute is spent on search.

This project builds \textbf{EvoScale}, a system that closes both gaps. It
runs hundreds of problems in parallel, generates $K$ candidate mutations
per iteration in batch, and uses the resulting reward signal to fine-tune
the LLM mutation operator online via Group Relative Policy Optimization
(GRPO)~\cite{deepseekr1}. The central thesis is that at sufficient scale,
\emph{inference is training}: every token generated in service of search
also produces a gradient step that improves subsequent search iterations.

% -----------------------------------------------------------------------
\section*{Method}

\subsection*{Distributed Async Evolution}

The standard evolutionary loop serializes three operations: sample a
parent, call the LLM, evaluate the result. EvoScale restructures all
three.

\textbf{Batch candidate generation.} Each iteration issues $K = 8$
parallel LLM requests to a vLLM server~\cite{vllm}, all conditioned on the
same parent and inspiration context. vLLM's continuous batching handles
these at near-zero overhead relative to a single request at GPU capacity.
Crucially, this batch is a natural group for GRPO: $K$ candidates with
shared context and individual rewards, requiring no value function or
importance weighting.

\textbf{Async evaluation pipeline.} LLM generation and program evaluation
are decoupled via a producer-consumer queue. While one batch of candidates
is being evaluated, the next batch of LLM calls is already running. This
eliminates the GPU idle time that dominates sequential execution and brings
utilization toward saturation.

\textbf{Multi-problem orchestration.} Up to $N = 200$ independent
evolution tasks run as Ray actors, each with its own population. All
actors share one vLLM endpoint. A central coordinator collects
(prompt, response, reward) tuples from all actors into a training
buffer and periodically triggers gradient updates. Running many
problems in parallel also ensures reward diversity across tasks,
which is essential for learning a general mutation operator rather
than a task-specific one.

Static system prompts that encode long problem descriptions (ARC-AGI-2
training examples can exceed 10{,}000 tokens) are placed at the front of
every request so that vLLM's radix attention prefix cache amortizes the
prefill cost across all iterations of the same problem, substantially
reducing effective per-token cost.

\subsection*{Online RL: Fine-Tuning the Mutation Operator}

The policy $\pi_\theta$ is Qwen2.5-Coder-32B~\cite{qwen25coder} with a
LoRA~\cite{lora} adapter. The state is the full evolution prompt. The
action is the proposed code mutation. The reward is normalized score
improvement:
\[
  r = \frac{\text{score}(\text{child}) - \text{score}(\text{parent})}
           {\text{score}_{\max}}
\]
with a fixed penalty for parse failures that waste evaluation budget.

GRPO uses the $K$ within-batch candidates as a comparison group. Their
relative rewards define advantages directly, with no critic, no
value function, and no off-policy correction needed. Every $M$ tuples
accumulated in the buffer, one gradient step updates the LoRA adapter.
Updated weights are hot-reloaded into the vLLM replicas via the vLLM
weight-update API using a rolling replacement protocol. LoRA restricts
parameter changes to approximately 0.5\% of the model, preserving
general coding ability while allowing the adapter to specialize on
evolutionary mutation patterns.

% -----------------------------------------------------------------------
\section*{Evaluation Plan}

\textbf{Benchmarks.} Three domains cover qualitatively different reward
structures: ARC-AGI-2 (120 evaluation tasks, sparse binary reward),
Triton GPU kernel optimization (dense continuous reward based on runtime
speedup), and AtCoder Heuristic Contest problems (competitive objective
function scores). This breadth tests whether the RL signal generalizes
across task types.

\textbf{Baselines.} Four conditions form a clean ablation: (1) sequential
AdaEvolve with GPT-5, the published state of the art; (2) sequential
AdaEvolve with Qwen2.5-Coder-32B, isolating the model gap; (3) parallel
EvoScale with Qwen2.5-Coder-32B and no fine-tuning, isolating the systems
contribution; and (4) full EvoScale with online GRPO.

\textbf{Metrics.} We report final score per benchmark, wall-clock time to
reach a fixed score threshold, and quality per GPU-hour. The key scientific
question is a \textbf{generalization test}: fine-tune the LoRA adapter
exclusively on ARC-AGI-2 and evaluate it without modification on GPU
kernel and competitive programming benchmarks. Positive transfer would
demonstrate that the adapter has learned a general skill of proposing good
code improvements, rather than a benchmark-specific pattern. This is the
result that would distinguish EvoScale from a simple engineering speedup.

% -----------------------------------------------------------------------
\section*{Risks}

\textbf{Distribution shift.} As the policy improves, the prompt
distribution it encounters shifts, making earlier tuples off-policy.
We address this with a recency-weighted replay buffer that down-weights
stale experience and a KL penalty against the base model in the GRPO
objective to prevent the adapter from drifting far from the base
distribution.

\textbf{Reward sparsity.} Most mutations do not improve the score. With
$N = 100$ parallel problems and $K = 8$ candidates per iteration, each
training batch contains approximately 800 evaluated mutations. At an
expected improvement rate of 10 to 15\%, this yields 80 to 120 positive
reward signals per batch, which is comparable to signal density in
successful RL applications to code generation~\cite{deepseekr1}. If
signals remain insufficient on a given benchmark, we add parse correctness
as a dense auxiliary reward to bootstrap early training.

% -----------------------------------------------------------------------
\section*{Potential Impact}

\textbf{A new compute scaling regime.} Every existing inference-time
scaling method (chain-of-thought, best-of-N, tree search) treats the
model as fixed. Compute spent on search is disposable: it improves the
current output and nothing else. EvoScale introduces a feedback loop
in which search performance and model quality co-improve. A successful
result establishes a compound scaling law: more iterations improve the
model, which improves subsequent iterations. This is categorically
different from any fixed-policy scaling curve and directly relevant
to the question of how to allocate large compute budgets efficiently.

\textbf{Closing the OSS gap to frontier models.} AdaEvolve achieves its
published results with GPT-5 and Gemini-3-Pro. Open-source models start
behind on these benchmarks. Online fine-tuning on evolutionary search
trajectories gives the adapted model a structural advantage: it has been
trained on thousands of (parent program, mutation, outcome) examples
from the exact task distribution it will be evaluated on, a form of
task-specific supervision that no general-purpose frontier model has
received. If EvoScale with a fine-tuned Qwen model matches or exceeds
AdaEvolve with GPT-5 on any benchmark, it demonstrates that targeted
online RL on accessible hardware can close the gap to proprietary models,
with concrete implications for AI accessibility.

\textbf{Compute efficiency.} The systems design ensures that no component
of the allocated budget is idle. The async pipeline eliminates GPU idle
time during evaluation. Prefix caching eliminates redundant prefill. Batch
generation amortizes request overhead across $K$ candidates. RL training
executes during evaluation windows that would otherwise go unused. The
result is that every GPU-hour serves double duty: advancing the search and
training the model. We will quantify this directly as quality gained per
GPU-hour across all baselines.

\textbf{A reusable systems artifact.} There is no existing open system that
co-schedules vLLM inference with online LoRA fine-tuning under live serving
conditions. The infrastructure we build, encompassing Ray-based
multi-problem orchestration, async evaluation pipelines, within-batch GRPO,
and live weight updates, applies directly to any deployment where model
behavior generates its own fine-tuning signal. This includes tool-using
agents, long-horizon planners, and any setting where an LLM is evaluated
against an objective rather than a human preference. The systems
contribution stands independently of the RL generalization result.

% -----------------------------------------------------------------------
\bibliographystyle{plain}
\begin{thebibliography}{9}

\bibitem{adaevolve2025}
AdaEvolve: Adaptive LLM-Driven Zeroth-Order Optimization.
\textit{arXiv:2602.20133}, 2025.
\url{https://arxiv.org/abs/2602.20133}

\bibitem{deepseekr1}
DeepSeek-AI.
DeepSeek-R1: Incentivizing Reasoning Capability in LLMs via Reinforcement
Learning.
\textit{arXiv:2501.12948}, 2025.

\bibitem{vllm}
W.~Kwon et al.
Efficient Memory Management for Large Language Model Serving with
PagedAttention.
\textit{SOSP}, 2023.

\bibitem{qwen25coder}
Hui et al.
Qwen2.5-Coder Technical Report.
\textit{arXiv:2409.12186}, 2024.

\bibitem{lora}
E.~J.~Hu et al.
LoRA: Low-Rank Adaptation of Large Language Models.
\textit{ICLR}, 2022.

\end{thebibliography}


\end{document}